\section{Введение}

\setcounter{page}{1}

Обучение с подкреплением (RL) — это универсальный подход для решения
различных задачи с оптимальной стратегией.
Благодаря стабильной среде и системе явных вознаграждений для агента-исполнителя, этот
подход способен обучить агента максимизировать сбор вознаграждений,
совершая наиболее полезные действия с точки зрения окружения.

Практический пример применения обучения с подкреплением – задача 
принятия решений совместно с прогнозированией временных рядов, широко встречаемая в настоящее время 
в различных отраслях.
Такие финансовые задачи имеют четкую систему вознаграждений, легко интерпретируемая как людьми, так и для обучающих
алгоритмов, что делает их привлекательными для решения RL-подходами \cite{choi2009reinforcement, madureira2023ensemble}.

Однако в реальных задачах временных рядов основная сложность проявляется в "шоках" - явлениях
непредсказумого изменения во временном ряду,
которые в основном вызваны значимыми мировыми событиями с большим макроэкономическим влиянием или наоборот - чистым 
совпадением нескольких событий с меньшим влиянием \cite{balke1994large}.
В результате даже малейшее непредвиденное изменение в динамике временных рядов или статических 
(не зависимых от действий агентов) аспектах RL-среды,
может привести к ошибочным предсказаниям моделями.

Эти события можно предотвратить, внося аналогичные синтетические изменения в сценарии обучения, 
однако это усложняет процесс обучения RL и требует более дорогой архитектуры.

Подходы глубокого обучения с подкреплением (далее, DRL) могли бы иметь еще большее преимущество в решении 
экономических проблем с многочисленными состояниями, вбирая в себя не полную информацию об обучении, а 
лишь наиболее значимые для наград паттерны окружения. 
Однако DRL-модели также трудно защитить от переобучения и, следовательно, они не устойчивы к потрясениям в 
рамках обучения стандартным окружениям.
Решением может стать дополнительное обучение агента в новых условиях внешней среды «на ходу»,
с "теплым" стартом по параметрам модели, предварительно обученной на условиях исходной среды
Этот метод потенциально может быть быстрее, чем полное переобучение начиная со
случайных параметров и дешевле, чем обучение прогнозированию
дополнительных возможных «шоков».

В этой статье мы рассматриваем набор подходов DRL, основанных на методе обучения Deep-Q Network,
с возможностью дополнительного обучения, позволяющее реагировать на возникшие "шоки" в среде
путем проведения «горячего» переобучения.
Также мы рассмотрим интерактивную задачу о целенаправленном движении объектов.
с набором ограничений, аналогичных средам на базе Atari.
Представленные в работе сценарии «шоков» среды затронут важнейшие части карты окружения, 
а также оставят доступные пути успешно решения задачи.